
\subsection{Optimized Peeling for $r=2$}
\label{sec:opt_r2}

For $r=2$, the baseline algorithm (Algorithm~\ref{alg:removeEdge}) handles three cases:
hold--hold edges (path dies), hold--pivot edges (pivot removed), and pivot--pivot edges (path splits via residual-graph search).
We optimize the latter two cases with closed-form formulas that eliminate most residual-graph searches.

\stitle{Key Idea: Closed-Form Deltas for Case~C.}
When pivots are removed from a surviving path (hold--pivot edges or additional pivot removals), the support change for every remaining edge can be computed analytically, analogous to the $r=1$ case.

\begin{theorem}[Closed-Form Support Delta for Pivot Removal]\label{thm:r2_pivot_delta}
Let $\TPath = (V_h, V_p)$ be a clique path with $|V_p| = P$, $|V_h| = K$, and $\mathit{np} = s - K$.
After removing $d$ pivots from $\TPath$ (path survives: $\mathit{np} \le P - d$), the support change for each surviving edge $e$ is determined by its type:
\begin{itemize}[leftmargin=*]
\item \textbf{Hold--hold edge} ($e \subseteq V_h$): \quad $\Delta(e) = \binom{P}{np} - \binom{P-d}{np}$;
\item \textbf{Pivot--pivot edge} ($e \subseteq V_p$, both survive): \quad $\Delta(e) = \binom{P-2}{np-2} - \binom{P-d-2}{np-2}$;
\item \textbf{Hold--pivot edge} ($|e \cap V_h|=1$, pivot survives): \quad $\Delta(e) = \binom{P-1}{np-1} - \binom{P-d-1}{np-1}$;
\item \textbf{Edge involving a removed pivot}: \quad $\Delta(e)$ equals the full original weight (edge loses all support from this path).
\end{itemize}
Each delta is an $O(1)$ binomial-coefficient lookup.
\end{theorem}

\begin{proof}
By Lemma~\ref{lem:clique_count}, the support contribution of an edge $e = \{u,v\}$ on path $\TPath$ depends on whether $u$ and $v$ are holds or pivots.
For a hold--hold edge, both vertices appear in every $s$-clique on $\TPath$, so the contribution equals $\binom{P}{np}$; after removing $d$ pivots this becomes $\binom{P-d}{np}$.
The other cases follow by fixing one or two pivots in the binomial coefficient.
\end{proof}


\stitle{Key Idea: Analytical Single-Edge Split for Case~B.}
When exactly one pivot--pivot edge $\{u,v\}$ is removed from a path (the most common Case~B scenario), we can construct the two resulting subpaths \emph{directly} from Theorem~\ref{thm:remove_edge}(3) without invoking the residual-graph search of Algorithm~\ref{alg:removeEdge}.

\begin{theorem}[Analytical Single-Edge Split]\label{thm:r2_single_split}
Let $\TPath = (V_h, V_p)$ be a clique path and let $\{u, v\} \subseteq V_p$ be a single removed pivot--pivot edge.
The valid $s$-cliques (those not containing both $u$ and $v$) partition into two disjoint groups, each representable as a clique path:
\begin{enumerate}[leftmargin=*]
\item $\TPath_1 = (V_h,\; V_p \setminus \{v\})$: \quad $s$-cliques \textbf{not} containing $v$;
\item $\TPath_2 = (V_h \cup \{v\},\; V_p \setminus \{u, v\})$: \quad $s$-cliques containing $v$ but \textbf{not} $u$.
\end{enumerate}
The support delta for each surviving edge on each subpath is computed via the nCr formula (Lemma~\ref{lem:clique_count}) in $O(1)$ per edge.
\end{theorem}

\begin{proof}
Every valid $s$-clique either does not contain $v$ (belongs to $\TPath_1$)
or contains $v$ but not $u$ (belongs to $\TPath_2$).
These two sets are disjoint ($\TPath_1$ excludes $v$; $\TPath_2$ includes $v$)
and exhaustive (invalid cliques are exactly those containing both $u$ and $v$).
In $\TPath_2$, vertex $v$ is moved from $V_p$ to $V_h$ because every $s$-clique in this group must contain $v$.
\end{proof}

\stitle{Effectiveness.}
In practice, the single-edge case (Theorem~\ref{thm:r2_single_split}) accounts for 93--97\% of all pivot--pivot removals at small $s$ values.
Combined with the closed-form pivot-removal delta (Theorem~\ref{thm:r2_pivot_delta}),
these optimizations eliminate the residual-graph search in the vast majority of peeling steps,
reducing the per-step cost from $O(3^{d/3} \cdot |\TPath|^2)$ (residual-graph BK) to $O(|\TPath|)$ (closed-form).

\stitle{Algorithm.}
Algorithm~\ref{alg:removeEdge} is extended as follows:
after classifying each affected path into Case~A (hold--hold $\to$ path dies), Case~C (hold--pivot or pivot-only removal $\to$ closed-form delta via Theorem~\ref{thm:r2_pivot_delta}), or Case~B (pivot--pivot edges remain):
\begin{enumerate}[leftmargin=*]
\item If $|E_{\mathrm{rm}}| = 0$ after removing hold--pivot edges $\to$ output the surviving path directly.
\item If $|E_{\mathrm{rm}}| = 1$ (single pivot--pivot edge) $\to$ apply Theorem~\ref{thm:r2_single_split} to produce two subpaths analytically.
\item If $|E_{\mathrm{rm}}| \ge 2$ $\to$ fall back to the residual-graph search (Algorithm~\ref{alg:removeEdge}).
\end{enumerate}

\begin{algorithm}[t]
\small
    \small
\caption{\textsc{$(2,s)$-CND}}
\label{alg:peeling_r2}
\KwIn{Succinct Clique Tree (\sct) $\tree$, Clique size $s$}
\KwOut{Core value $\kappa_s(v)$ for every vertex $v$}
\SetKwFunction{RemoveNode}{RemoveNode}
\SetKwFunction{CountPreEdges}{$s$-CliqueCount}
\SetKwProg{Fn}{Function}{:}{}

$count\gets \CountPreEdges(\tree,2,s)$\;

build min heap $H_{\min}$ using $count$\;

$c_{min} \gets 0$\;

\While{$H_{\min}$ is not empty}{
    $e_{min} \gets \text{ExtractMin}(H_{\min})$\;
    $c_{min} \gets \max(c_{min}, count(e_{min}))$\;
    $\kappa_s(e_{min}) \gets c_{min}$\;
    \ForEach{$e \in e_{min}$}{
        % remove edge base on lemma~\ref{lem:remove_edge}
        \ForEach{leaf $\TPath$ containing $e$}{
            \text{Remove Edge $e$ from $\TPath$ based on Theorem~\ref{thm:remove_edge}}\;
        }
    }
    \ForEach{leaf $\TPath$ containing $e \in e_{min}$}{
        Update $count(e')$ and $H_{\min}$ for all edges $e' \in \TPath$
    }  

}

\Return{$\kappa_s(e)$ for all edges $e$}\;

\end{algorithm}
